\section{Experiments}
\label{sec:experiments}

\subsection{Setup}
\label{sec:setup}
We use a ViT-B/16 backbone pre-trained on ImageNet-1k, frozen throughout, with LoRA
adapters (rank $r=16$, $\alpha=16$) attached to the MLP down-projection of every
transformer block (12 adapters in total). The classifier head is a single linear layer
grown as new classes arrive. All methods share the identical backbone and head; they
differ only in the adapter regularization.

\textbf{Datasets.} Split CIFAR-100 (20 tasks $\times$ 5 classes) and Split ImageNet-R
(20 tasks $\times$ 10 classes), both under the class-incremental protocol (task identity
unavailable at test time).

\textbf{Baselines.} (i) \emph{Seq-LoRA}: sequential fine-tuning of the shared adapter
with no protection (lower bound); (ii) \emph{EWC-LoRA}: diagonal Fisher on the LoRA
parameters~\cite{kirkpatrick2017overcoming}; (iii) \emph{O-LoRA}: per-task adapters with
orthogonal initialization~\cite{wang2023olora}; (iv) \emph{L2P}: a shared pool of prompts
selected per input by key--query matching~\cite{wang2022l2p}; (v) \emph{CODA-Prompt}: a set
of prompt components combined by attention~\cite{smith2023coda}; (vi) \emph{FOLoRA} (ours).

\textbf{Training.} AdamW with learning rate $10^{-3}$, batch size 32, and bf16 mixed
precision. We train the adapter-based methods (Seq-LoRA, EWC-LoRA, O-LoRA, FOLoRA) for
$5$ epochs per task, which avoids over-fitting and keeps the empirical Fisher meaningful
(Sec.~\ref{sec:discussion}); the prompt-based methods (L2P, CODA-Prompt) are trained for
$20$ epochs per task, their optimal setting since their key-based prompt selection
benefits from more epochs. For FOLoRA we select $\lambda$ and $k$ by a small grid search
over $\lambda\in\{10,30,100,300,1000\}$ and $k\in\{16,64\}$, taking $\lambda=300$,
$k=16$ as the default; the importance kernel is estimated from up to 300 per-sample
gradients per task. EWC's regularization strength is likewise tuned on CIFAR-100, reaching
its optimum at $\lambda=30$. We report the mean and standard deviation over 10 seeds on
CIFAR-100 (5 seeds on ImageNet-R) for the adapter-based methods, and 3 seeds for the
prompt-based methods.

\textbf{Metrics.} Average accuracy over all tasks after the final task
($\overline{\mathrm{ACC}}$), average forgetting ($\mathrm{FGT}$), and average
incremental accuracy.

\subsection{Main results}
\label{sec:main_results}

\begin{table}[t]
\centering
\caption{Class-incremental results (mean $\pm$ std). Adapter-based methods are trained for
5 epochs (10 seeds on CIFAR-100, 5 on ImageNet-R); prompt-based methods for 20 epochs
(3 seeds). $\overline{\mathrm{ACC}}$: average accuracy; $\mathrm{FGT}$: average
forgetting. Lower $\mathrm{FGT}$ is better. Best in \textbf{bold}.}
\label{tab:main}
\begin{tabular}{lcccc}
\toprule
& \multicolumn{2}{c}{Split CIFAR-100} & \multicolumn{2}{c}{Split ImageNet-R} \\
\cmidrule(lr){2-3} \cmidrule(lr){4-5}
Method & $\overline{\mathrm{ACC}}$ & $\mathrm{FGT}$ & $\overline{\mathrm{ACC}}$ & $\mathrm{FGT}$ \\
\midrule
Seq-LoRA & 7.19$\pm$0.93 & 90.90$\pm$1.05 & 6.55$\pm$0.36 & 84.42$\pm$0.26 \\
EWC-LoRA & 9.17$\pm$0.89 & 88.76$\pm$1.00 & 8.17$\pm$0.75 & 82.92$\pm$0.95 \\
O-LoRA   & 7.74$\pm$0.79 & 90.01$\pm$0.96 & 6.38$\pm$0.43 & 83.44$\pm$0.27 \\
L2P      & 8.67$\pm$0.55 & 89.08$\pm$0.58 & 7.05$\pm$0.32 & 80.22$\pm$0.36 \\
CODA-Prompt & 9.35$\pm$0.52 & 88.52$\pm$0.65 & 6.93$\pm$0.27 & 79.54$\pm$0.59 \\
FOLoRA (ours) & \textbf{10.07}$\pm$1.54 & \textbf{87.65}$\pm$1.43 & \textbf{9.56}$\pm$0.83 & \textbf{79.44}$\pm$0.62 \\
\bottomrule
\end{tabular}
\end{table}

On both benchmarks FOLoRA attains the highest average accuracy and the lowest average
forgetting, including over the prompt-based methods at their optimal epoch budget. The
improvement is largest on ImageNet-R, where FOLoRA improves accuracy by $+1.39$ points
over the strongest baseline (EWC-LoRA) while reducing forgetting by $3.48$ points; on
CIFAR-100 the gains over EWC-LoRA are $+0.90$ and $-1.11$ points, respectively, and over
the strongest prompt baseline (CODA-Prompt) $+0.72$ and $-0.87$. The forgetting reduction
is statistically significant on both benchmarks (Welch's $t$-test: $p<0.001$ on ImageNet-R
and $p=0.044$ on CIFAR-100); the accuracy improvement is significant on ImageNet-R
($p=0.005$) and a clear mean advantage on CIFAR-100 ($p=0.110$), where the high variance
of the frozen-backbone, no-replay regime limits single-seed significance
(Sec.~\ref{sec:discussion}). We further tuned EWC's regularization strength on CIFAR-100 by
grid search ($\lambda\in\{1,10,30,100,300,1000\}$); its optimum $\lambda=30$ reaches
$9.58\pm0.78$, still $0.49$ points below FOLoRA, confirming that the advantage is not an
artifact of an under-tuned baseline. All methods share a fixed frozen backbone and a
single classifier head; the parameter count of FOLoRA, Seq-LoRA, EWC-LoRA, L2P and
CODA-Prompt stays fixed across tasks, whereas O-LoRA allocates a new adapter per task and
thus grows linearly with the number of tasks.

\subsection{Ablation study}
\label{sec:ablation}

We ablate the two hyperparameters of FOLoRA --- the regularization strength $\lambda$ and
the number of protected directions $k$ --- on Split CIFAR-100 (seed 0).
Table~\ref{tab:ablation} reports average accuracy and forgetting.

\begin{table}[t]
\centering
\caption{Ablation of $\lambda$ (with $k=16$) and $k$ (with $\lambda=10$) on Split
CIFAR-100, seed 0.}
\label{tab:ablation}
\begin{tabular}{cccc}
\toprule
$\lambda$ & $\overline{\mathrm{ACC}}$ & $\mathrm{FGT}$ & \\
\midrule
10  & 9.20 & 89.22 & \\
30  & 9.80 & 88.53 & \\
100 & 10.94 & 87.19 & \\
300 & \textbf{11.64} & \textbf{86.12} & \\
1000 & 10.69 & 87.09 & \\
\midrule
$k$ & $\overline{\mathrm{ACC}}$ & $\mathrm{FGT}$ & \\
\midrule
16 & \textbf{9.20} & \textbf{89.22} & \\
64 & 8.61 & 89.54 & \\
\bottomrule
\end{tabular}
\end{table}

Accuracy is highest at $\lambda=300$ and degrades on both sides, confirming that a
moderately strong regularization best balances stability and plasticity: too small a
$\lambda$ under-protects the important directions, while too large a $\lambda$
over-constrains the adapter and hurts plasticity. Increasing $k$ from 16 to 64 does not
help and slightly hurts, indicating that a compact protected subspace ($k=16$) already
captures the dominant forgetting directions.

\subsection{Limitations and discussion}
\label{sec:discussion}

The forgetting reduction of FOLoRA is statistically significant on both benchmarks. On
CIFAR-100 the accuracy improvement is consistent in mean but noisy across seeds: the
average accuracy of FOLoRA spans $0.086$--$0.127$ over the ten seeds, whereas the
baselines show a tighter spread. We attribute this variance to two sources. First, the
importance kernel is estimated from a finite per-sample budget (up to 300 gradients per
task), and its top-$k$ approximation is sensitive to the particular samples that land in
the tails of the estimate. Second, under a frozen backbone, a fixed-rank ($r{=}16$)
adapter and no replay, the class ordering exerts a strong effect on the final subspace, so
individual seeds drift apart. We further found that the epoch budget interacts with the
method family: adapter-based Fisher methods degrade under longer training because
over-fitting drives the per-sample gradients toward zero and thereby degenerates the
estimated Fisher, whereas prompt-based methods benefit from longer training of their
prompt keys. This motivates the per-family epoch budgets in Sec.~\ref{sec:setup}. Even in
this unfavorable regime FOLoRA never falls below the strongest baseline on either
benchmark, which supports its robustness. A natural direction for future work is a
stochastic (Monte-Carlo) Fisher estimate and class-order averaging to tighten the
variance.
